% =====================================================================
%  Lecture 15 -- Software-Defined Networking (SDN)        (English edition)
% =====================================================================
\documentclass[11pt,a4paper]{article}
\providecommand\ACNroot{../..}
\input{\ACNroot/templates/lecture-en}

\ACNset{number}{15}
\ACNset{title}{Software-Defined Networking (SDN)}
\ACNset{shorttitle}{SDN}
\ACNset{titleother}{شبکه‌سازی نرم‌افزارمحور (SDN)}
\ACNset{part}{Part III -- Software-driven architectures (Lectures 15--18)}
\ACNset{time}{6--8 hours (3 h reading, 3--5 h exercises)}
\ACNset{level}{Intermediate}
\ACNset{prereq}{IP routing and switching; Lecture 13 (CUPS); M/M/1 queue (Lecture 12)}
\ACNset{gen}{5}

\begin{document}
\ACNcover

% ---------------------------------------------------------------------
\section{Learning objectives}
After studying this lecture you should be able to:
\begin{objectives}
  \item \textbf{Explain} the separation of control and data planes and the motivation for logically
        centralised control.
  \item \textbf{Describe} the SDN layers and interfaces (northbound, southbound) and the OpenFlow
        switch model: flow tables, pipeline, groups and meters.
  \item \textbf{Trace} reactive and proactive flow installation and \textbf{compute} their latency and
        table requirements.
  \item \textbf{Explain} data-plane programmability with P4.
  \item \textbf{Analyse} controller placement and controller load.
  \item \textbf{Identify} SDN principles in mobile networks (CUPS, PFCP, transport and slicing).
\end{objectives}

% ---------------------------------------------------------------------
\section{Prerequisite knowledge}
You should know how IP routers and Ethernet switches forward packets (longest-prefix match, MAC learning)
and that routing protocols such as OSPF and BGP compute forwarding tables in a distributed way. From the
course, you need CUPS and PFCP (Lecture~13) and the M/M/1 delay formula (Lecture~12).

\paragraph{Key terms} The vocabulary of this lecture with its Persian equivalents:
\begin{keyterms}
Control plane & \fa{صفحهٔ کنترل} & Logic that decides how traffic is forwarded. \\
Data plane & \fa{صفحهٔ داده} & Per-packet forwarding according to installed state. \\
SDN controller & \fa{کنترل‌گر SDN} & Logically centralised software that programs the data plane. \\
Southbound API & \fa{API جنوبی} & Controller--device interface (OpenFlow, P4Runtime, NETCONF). \\
Northbound API & \fa{API شمالی} & Controller--application interface (e.g.\ REST). \\
Flow table & \fa{جدول جریان} & Prioritised match--action rules in a switch. \\
Packet-In / Flow-Mod & \fa{پیام‌های OpenFlow} & Switch-to-controller packet report / rule installation. \\
Reactive / proactive & \fa{واکنشی / پیش‌کنشی} & Rules installed on first packet / in advance. \\
P4 & \fa{زبان P4} & Language for programming protocol-independent packet processors. \\
TCAM & \fa{حافظهٔ آدرس‌پذیر محتوایی سه‌حالته} & Hardware memory for wildcard matching. \\
\end{keyterms}

% ---------------------------------------------------------------------
\section{Introduction and motivation}
In a traditional network every router runs its own control software: it exchanges routing messages with
its neighbours, computes paths and fills its forwarding table. Control and forwarding are bundled in each
box, and each vendor implements the protocols in its own closed software. This design is robust, but it
makes the network hard to change: a new forwarding behaviour needs a new protocol, a standard and years
of implementation, and network-wide policies (traffic engineering, isolation, access control) must be
expressed indirectly through many per-box configurations.

Software-defined networking (SDN) separates the two planes \cite{mckeown2008openflow,kreutz2015sdn}.
Forwarding devices become simple, programmable match--action engines; a logically centralised controller
with a global view computes their state and installs it over an open interface. Network behaviour becomes
software that can be written, tested and deployed like other software. Mobile networks adopted the same
idea: CUPS and PFCP (Lecture~13) are SDN applied to packet gateways, and the 5G core with its UPF is built on
it. This lecture opens Part~III, which covers the software technologies -- SDN, NFV, service-based
architecture -- on which 5G networks rest.

% ---------------------------------------------------------------------
\section{Fundamental concepts}

\subsection{Three principles}
\begin{enumerate}
  \item \textbf{Separation of control and data planes}: devices forward; controllers decide.
  \item \textbf{Logically centralised control}: one (possibly distributed and replicated) controller
        has a global view of topology and state.
  \item \textbf{Programmability}: open interfaces let software applications define network behaviour.
\end{enumerate}

\subsection{Layers and interfaces}
RFC~7426 describes SDN in terms of planes -- forwarding, operational, control, management and application --
and the abstraction layers between them \cite{rfc7426}. In the common simplified view
(\cref{fig:arch}) applications use a \emph{northbound} API (often REST) of the controller, and the controller
programs devices through a \emph{southbound} API: OpenFlow \cite{onf2015openflow} for flow tables, P4Runtime
for P4-programmed devices, or NETCONF/YANG for configuration.

\begin{acnfigure}
  \resizebox{\linewidth}{!}{\acnfig{lecture-15/sdn-architecture}}
  \caption{SDN layers: applications, controller and data plane with northbound and southbound APIs.}
  \label{fig:arch}
\end{acnfigure}

\subsection{The OpenFlow switch}
\begin{definitionbox}{Flow entry}
A flow entry consists of match fields (ingress port and packet header fields from Ethernet to transport
layer), a priority, counters, instructions (apply or write actions, go to another table), idle and hard
timeouts, and an opaque cookie \cite{onf2015openflow}. A packet is matched against the entries of a table;
the highest-priority match wins.
\end{definitionbox}
A switch contains a pipeline of flow tables (\cref{fig:table}), a \emph{group table} (for multicast,
load balancing with \emph{select} groups and fast failover) and a \emph{meter table} (rate limiting). If no
entry matches (\emph{table miss}), the table-miss entry decides: drop, continue, or send the packet to the
controller. The switch talks to the controller over a TCP connection (port 6653, optionally TLS).

\begin{acnfigure}
  \resizebox{\linewidth}{!}{\acnfig{lecture-15/flow-table}}
  \caption{OpenFlow pipeline with multiple flow tables, and the fields of a flow entry.}
  \label{fig:table}
\end{acnfigure}

% ---------------------------------------------------------------------
\section{Detailed technical explanation}

\subsection{OpenFlow messages}
\begin{itemize}
  \item \textbf{Controller-to-switch}: Features (capabilities), Flow-Mod (add, modify, delete entries),
        Group-Mod, Meter-Mod, Packet-Out (send a packet from a port), Multipart/Stats (read counters), Barrier
        (ensure earlier messages are processed).
  \item \textbf{Asynchronous}: Packet-In (a packet sent to the controller, by default with its first 128 bytes
        and a buffer identifier), Flow-Removed (an entry expired), Port-Status (a port changed).
  \item \textbf{Symmetric}: Hello, Echo (liveness), Error.
\end{itemize}

\subsection{Reactive and proactive control}
In \emph{reactive} mode the first packet of a new flow misses, the switch sends a Packet-In, and the
controller computes a path and installs entries along it before releasing the packet with Packet-Out
(\cref{fig:reactive}). Fine-grained control is possible, but every new flow pays a controller round trip,
and the controller load grows with the flow arrival rate. In \emph{proactive} mode the controller installs
(usually wildcard) rules in advance, from its knowledge of topology and policy; no packet waits, the
controller is off the critical path, but rules must be aggregated to fit the tables. Practical deployments
mix both: proactive rules for the bulk of traffic, reactive handling for exceptions.

\begin{acnfigure}
  \resizebox{\linewidth}{!}{\acnfig{lecture-15/reactive-flow}}
  \caption{Reactive flow set-up: table miss, Packet-In, Flow-Mod along the path, Packet-Out.}
  \label{fig:reactive}
\end{acnfigure}

\subsection{Data-plane programmability: P4}
OpenFlow fixes the header fields a switch can match; adding a new protocol requires a new version of the
specification and of the hardware. P4 \cite{bosshart2014p4} goes one step further: the program defines the
headers, the parser state machine, the match--action tables and the control flow, and a compiler maps it to a
programmable target (switch ASIC, smart NIC, software switch) organised as a parser, ingress and egress
match--action pipelines and a deparser (\cref{fig:p4}). The controller then fills the tables at run time
through P4Runtime. P4 has been used, for example, to implement GTP-U encapsulation and UPF functions directly
in switch hardware.

\begin{acnfigure}
  \resizebox{\linewidth}{!}{\acnfig{lecture-15/p4-pipeline}}
  \caption{Protocol-independent switch architecture programmed with P4.}
  \label{fig:p4}
\end{acnfigure}

\subsection{Distributed controllers and consistency}
A single controller is a single point of failure and a bottleneck. Production controllers run as clusters
that replicate state with consensus protocols and partition switches among instances; a switch connects to
several controllers with master/slave (later primary/secondary) roles. Updating many switches is not atomic:
during an update some packets may see old rules on one switch and new ones on the next, causing loops or
policy violations. \emph{Consistent update} techniques tag packets with a configuration version at the ingress
and install the new version on all switches before flipping the tags, so that each packet sees either the old
or the new configuration, never a mix.

\subsection{SDN in mobile networks}
\begin{itemize}
  \item \textbf{Gateways}: CUPS (Lecture~13) separates SGW/PGW control from user-plane functions; PFCP plays the
        role of a southbound protocol with rules (PDR, FAR, QER, URR) that resemble OpenFlow entries.
  \item \textbf{5G core}: the SMF controls UPFs over N4 with PFCP (Lecture~23).
  \item \textbf{Transport}: SDN controllers configure backhaul and fronthaul paths, segment routing and transport
        slices.
  \item \textbf{RAN}: the O-RAN architecture opens RAN control to software applications through a RAN intelligent
        controller -- the SDN idea applied to radio resource management.
\end{itemize}

% ---------------------------------------------------------------------
\section{Mathematical models}

\subsection{Reactive flow-set-up latency}\label{sec:setup}
\paragraph{Model} The first packet of a flow waits for the switch--controller round trip $\tau$ and the
controller's processing. With $S$ switches each generating $\lambda_s$ new flows per second and a controller
modelled as an M/M/1 server of rate $\mu$,
\begin{equation}\label{eq:setup}
  T_{\mathrm{setup}} \approx \tau + \frac{1}{\mu - S\lambda_s}, \qquad S\lambda_s < \mu .
\end{equation}
\paragraph{Example} $\tau = 2$\,ms, $S = 100$, $\lambda_s = 1000$\,s$^{-1}$, $\mu = 1.5\times10^5$\,s$^{-1}$:
$T_{\mathrm{setup}} = 2 + 0.02 = 2.02$\,ms. At $S\lambda_s = 1.49\times10^5$ the queueing term rises to 1\,ms
(\cref{fig:setup}).
\paragraph{Interpretation} For moderate load the propagation delay to the controller dominates, which argues for
controllers close to the switches; near saturation the queueing term explodes and Packet-In storms (e.g.\ during
a scan or a failure) can overload the controller -- one reason to prefer proactive rules.

\begin{acnfigure}
  \acnfig{lecture-15/setup-delay}
  \caption{Controller queueing delay versus Packet-In rate (\cref{eq:setup}).}
  \label{fig:setup}
\end{acnfigure}

\subsection{Flow-table occupancy}\label{sec:little}
\paragraph{Model} With exact-match (per-flow) rules, new flows arrive at rate $\lambda$ and each entry stays
for the flow duration plus the idle timeout, $\bar T$ on average. By Little's law the mean number of entries is
\begin{equation}\label{eq:little}
  \bar N = \lambda\,\bar T .
\end{equation}
\paragraph{Example} $\lambda = 2000$\,s$^{-1}$, mean flow duration 10\,s, idle timeout 10\,s: $\bar N = 40\,000$
entries -- far more than the few thousand TCAM entries of many switches. Per-flow reactive control does not
scale to such switches; wildcard rules, aggregation, or larger exact-match tables (hash-based) are needed.

\subsection{Controller placement}\label{sec:place}
\paragraph{Model} For switches $v\in V$ and controller locations $C$, with $d(v,c)$ the latency between them,
the average and worst-case switch--controller latencies are
\begin{equation}\label{eq:place}
  \bar L(C) = \frac{1}{|V|}\sum_{v\in V}\min_{c\in C} d(v,c), \qquad
  L_{\max}(C) = \max_{v\in V}\min_{c\in C} d(v,c).
\end{equation}
Minimising $\bar L$ for $|C| = k$ is the $k$-median problem; minimising $L_{\max}$ is the $k$-centre problem.
\paragraph{Example} Five switches on a line, 1\,ms apart: one controller at the centre gives $\bar L = 1.2$\,ms
($L_{\max} = 2$\,ms); at an end, $\bar L = 2$\,ms. Two controllers at the second and fourth switches give
$\bar L = 0.6$\,ms and $L_{\max} = 1$\,ms.

% ---------------------------------------------------------------------
\section{Network architecture and diagrams}
In a data centre, the SDN controller programs the leaf--spine fabric and the virtual switches in the
servers; in a WAN, it computes traffic-engineered paths between sites; in a mobile operator, SDN controllers
manage transport networks and, through the core's control functions, the user-plane nodes. The controller is
itself a distributed application deployed in redundant instances, often on the virtualised infrastructure
of Lecture~16.

% ---------------------------------------------------------------------
\section{Step-by-step procedures}

\begin{procedurebox}{Reactive flow installation}\label{proc:reactive}
\begin{steps}
  \item The first packet of a new flow arrives at switch 1 and misses all entries; the table-miss entry sends
        it to the controller as Packet-In (header and buffer identifier).
  \item The controller identifies source and destination, checks policy and computes a path with its topology
        view.
  \item It sends Flow-Mod messages to all switches on the path (destination side first, to avoid further
        misses), optionally followed by a Barrier.
  \item It sends Packet-Out to switch 1 to release the buffered packet along the new path.
  \item Subsequent packets match the new entries and are forwarded at line rate; entries expire after the idle
        timeout and Flow-Removed reports their counters.
\end{steps}
\end{procedurebox}

\begin{procedurebox}{Topology discovery}\label{proc:lldp}
\begin{steps}
  \item On connection, each switch reports its ports (Features reply).
  \item The controller sends, via Packet-Out, a link-discovery frame (LLDP) on every port of every switch.
  \item A neighbouring switch receives the frame and, by a pre-installed rule, returns it as Packet-In with its
        ingress port.
  \item From (sending switch, port) and (receiving switch, port) the controller learns each link; it repeats
        periodically and updates the topology on Port-Status events.
\end{steps}
\end{procedurebox}

% ---------------------------------------------------------------------
\section{Worked examples}

\begin{examplebox}{Set-up delay budget}\label{ex:setup}
\textbf{Problem.} A campus network has 100 switches; each sees 1000 new flows per second. The controller handles
$1.5\times10^5$ Packet-In/s and is 1\,ms away (one way). What delay does the first packet of a flow see, and what
fraction of short web flows (10 packets, 20\,ms) is it?

\textbf{Solution.} $\tau = 2$\,ms, queueing 0.02\,ms (\cref{eq:setup}): about 2\,ms, i.e.\ 10\,\% of the flow's duration.
The controller runs at 67\,\% load; a 50\,\% increase in flow rate would saturate it.
\end{examplebox}

\begin{examplebox}{Table size}\label{ex:table}
\textbf{Problem.} A top-of-rack switch with 4000 TCAM entries sees 2000 new flows/s, each lasting 10\,s on average,
idle timeout 10\,s. Can reactive per-flow rules be used? What if the idle timeout is 1\,s?

\textbf{Solution.} $\bar N = 2000\times20 = 40\,000 > 4000$: no. With 1\,s idle timeout, $\bar N = 22\,000$: still no.
Rules must be aggregated, e.g.\ per destination prefix (proactive), which needs entries per destination rather
than per flow.
\end{examplebox}

\begin{examplebox}{Placement}\label{ex:place}
\textbf{Problem.} Verify the values of the line example in \cref{sec:place}.

\textbf{Solution.} Centre (switch 3): distances 2, 1, 0, 1, 2; mean 1.2\,ms, max 2\,ms. End (switch 1): 0, 1, 2, 3, 4; mean
2\,ms, max 4\,ms. Controllers at 2 and 4: 1, 0, 1, 0, 1; mean 0.6\,ms, max 1\,ms.
\end{examplebox}

% ---------------------------------------------------------------------
\section{Technology comparison}
\begin{table}[htbp]
\centering\small
\caption{Traditional distributed control versus SDN.}\label{tab:cmp}
\rowcolors{2}{acnmist}{white}
\begin{tabularx}{\linewidth}{@{}l Y Y@{}}
\toprule
\rowcolor{white} & Distributed control (OSPF, BGP) & SDN \\
\midrule
Control location   & In every device                 & Logically centralised controller \\
View               & Local, built by protocol exchange & Global topology and state \\
New behaviour      & New protocol, standard, firmware & New controller application \\
Failure handling   & Local, fast, autonomous          & Needs controller or pre-installed backups \\
Scalability limit  & Protocol convergence, table size & Controller capacity, control-channel latency \\
Policy expression  & Indirect (metrics, filters per box) & Direct, network-wide \\
Vendor coupling    & Hardware and software bundled    & Open interfaces, merchant silicon \\
\bottomrule
\end{tabularx}
\end{table}

\paragraph{Not a replacement for routing} SDN does not abolish distributed protocols: many deployments keep BGP or
IS-IS for basic reachability and fast local failover, and use the controller for policy and optimisation on top
(hybrid SDN).

% ---------------------------------------------------------------------
\section{Key takeaways}
\begin{takeawaybox}
\begin{itemize}
  \item SDN separates control from forwarding, centralises control logically and makes the network programmable.
  \item OpenFlow switches: prioritised flow tables in a pipeline, groups, meters; Packet-In, Flow-Mod, Packet-Out.
  \item Reactive control costs $T_{\mathrm{setup}} \approx \tau + 1/(\mu - S\lambda_s)$ per new flow; proactive control
        needs aggregated rules.
  \item Table occupancy follows Little's law, $\bar N = \lambda\bar T$; per-flow rules rarely fit TCAMs.
  \item P4 makes the parser and pipeline programmable; controller placement is a $k$-median/$k$-centre problem.
  \item CUPS/PFCP and the 5G UPF apply SDN to the mobile core.
\end{itemize}
\end{takeawaybox}

% ---------------------------------------------------------------------
\section{Exercises}

\begin{exercisebox}{1}{Which plane?}\label{exr:planes}
Classify as control or data plane: (a) OSPF shortest-path computation; (b) TTL decrement; (c) installing a
Flow-Mod; (d) matching a packet against a TCAM; (e) LLDP-based topology discovery.
\end{exercisebox}

\begin{exercisebox}{1}{Reactive or proactive?}\label{exr:proactive}
For each case choose reactive or proactive control and justify: (a) a data-centre fabric with
destination-based forwarding; (b) a campus network where guests must authenticate before their first flow
is allowed; (c) a mobile backhaul network.
\end{exercisebox}

\begin{exercisebox}{2}{Control-channel bandwidth}\label{exr:packetin}
A controller receives $10^5$ Packet-In messages per second, each carrying 128 bytes of packet plus 32 bytes of
OpenFlow and TCP/IP overhead (amortised). Compute the control-channel bandwidth. How does it change if the
switch sends whole 1500-byte packets instead of buffering them?
\end{exercisebox}

\begin{exercisebox}{2}{Timeouts}\label{exr:little}
Using \cref{eq:little}, find the idle timeout that keeps the mean occupancy at 15\,000 entries for 1000 new
flows/s with mean duration 8\,s. What is the price of a very short idle timeout?
\end{exercisebox}

\begin{exercisebox}{3}{Placement on a ring}\label{exr:place}
Six switches form a ring with 1\,ms links. Compute $\bar L$ and $L_{\max}$ for one controller, and for two
controllers placed optimally. What happens to the worst case if one of the two controllers fails?
\end{exercisebox}

\begin{exercisebox}{3}{Consistent update}\label{exr:consistent}
A path change moves traffic from route A--B--D to A--C--D. Show how updating switch A before installing rules on C
can drop packets, and describe a version-tagging update that avoids both drops and mixing. How many rules does it
need temporarily?
\end{exercisebox}

\subsection*{Solutions}
\addcontentsline{toc}{subsection}{Solutions}

\begin{solutionbox}{\ref{exr:planes}}
(a) control; (b) data; (c) control (the act of installing; the entry is then used by the data plane); (d) data;
(e) control.
\end{solutionbox}

\begin{solutionbox}{\ref{exr:proactive}}
(a) Proactive: destinations are known, traffic volume is huge, and per-flow set-up delay is unacceptable.
(b) Reactive for the first flow of each guest (the controller must see it to apply authentication), then
proactive rules for authorised hosts. (c) Proactive: paths are planned and must survive controller outages.
\end{solutionbox}

\begin{solutionbox}{\ref{exr:packetin}}
$10^5\times160$\,bytes $= 16$\,MB/s $= 128$\,Mbit/s. With whole packets: $10^5\times1532\times8 = 1.23$\,Gbit/s -- ten times
more, and the controller must also store or forward them. Buffering in the switch saves control-channel capacity.
\end{solutionbox}

\begin{solutionbox}{\ref{exr:little}}
$15\,000 = 1000\,(8 + T_{\mathrm{idle}}) \Rightarrow T_{\mathrm{idle}} = 7$\,s. A very short idle timeout evicts entries of
flows that pause briefly; their next packet misses again and causes another Packet-In, increasing controller load
and latency.
\end{solutionbox}

\begin{solutionbox}{\ref{exr:place}}
One controller: distances 0, 1, 1, 2, 2, 3; $\bar L = 9/6 = 1.5$\,ms, $L_{\max} = 3$\,ms. Two opposite controllers: 0, 1, 1, 0,
1, 1; $\bar L = 4/6 = 0.67$\,ms, $L_{\max} = 1$\,ms. If one fails, the survivor serves all with $L_{\max} = 3$\,ms -- the
placement must be evaluated for failure cases too.
\end{solutionbox}

\begin{solutionbox}{\ref{exr:consistent}}
If A forwards to C before C has a rule for the flow, C misses and drops (or sends to the controller). Order matters:
install on C first, then switch A ("make before break"). For general updates: install the new rules on C and D
matching version tag 2 (while version-1 rules remain), then make A stamp tag 2 on new packets; once old packets
have drained, delete version-1 rules. Temporarily the affected switches hold both versions -- up to twice the
rules.
\end{solutionbox}

% ---------------------------------------------------------------------
\section{Further reading}
McKeown et al.\ \cite{mckeown2008openflow} introduced OpenFlow, whose current switch specification is published by the ONF
\cite{onf2015openflow}; Kreutz et al.\ \cite{kreutz2015sdn} survey SDN comprehensively; RFC~7426 \cite{rfc7426} defines
SDN layers and terminology; Bosshart et al.\ \cite{bosshart2014p4} present P4. For the mobile context see TS~23.214
\cite{ts23214} and TS~29.244 \cite{ts29244}.

\printbibliography[heading=subbibliography,title={\ACNstrRefs}]
\end{document}
